% !TEX root = ../main.tex
% !TEX program = lualatex
\documentclass[../main.tex]{subfiles}
\IfSubfilesClassLoaded{\externaldocument{\subfix{../build/main}}}{}
% =============================================================================
% APPENDIX: NAVY PAE CONSTRUCT MAPPING (ANALYST SUMMARY)
% DESCRIPTION: Working map of PAE portfolios against legacy SYSCOM/PEO seams.
% =============================================================================
\begin{document}
\IfSubfilesClassLoaded{\chapter{EDO Study Guide}}{}

\section{PAE Construct Mapping (Working Appendix)}
\label{app:pae-construct-mapping}
\textbf{Current as of 2026-04-27.}

\subsection{Executive Summary}
\begin{itemize}
  \item The Navy \ac{pae} construct is not a 1:1 relabeling of \acp{syscom} or \acp{peo}; it is a portfolio-layer authority model that cuts across traditional program, technical, contracting, sustainment, and industrial-base seams~\autocite{FY2026-NDAA,DON-PAE-Establishment-2026}.
  \item The six publicly announced baseline \acp{pae} are Robotics and Autonomous Systems, Industrial Operations, Marine Corps, Maritime, Strategic Systems Programs, and Undersea; the five non-RAS organizations were announced by \ac{don} on 16~Mar~2026 after the Dec~2025 RAS stand-up~\autocite{DON-PAE-Establishment-2026}.
  \item PAE Mission Systems should be treated as an internally announced, source-limited update: the 16~Mar~2026 \ac{don} release said mission-systems transition studies were still pending public establishment, while subsequent public artifacts identify a ``Portfolio Acquisition Executive Mission Systems (PAE MS)'' DVIDS unit page, list Mr.\ James ``Jim'' Day as a Portfolio Acquisition Executive speaker, and show PMW~760 integration work inside PAE Mission Systems~\autocite{DON-PAE-Establishment-2026,DVIDS-PAE-MS-Unit-2026,AFCEA-Navy-IW-Industry-Day-2026,PEOC4I-PMW760-NDIA-2026}.
  \item Key shift: \acp{syscom} increasingly function as technical authorities and infrastructure owners, \acp{peo} as execution arms, and \acp{pae} as portfolio owners with decision authority.
  \item The largest effects appear in cross-cutting portfolios (Undersea, Maritime, Industrial Operations, and Mission Systems) where multiple \acp{syscom}/\acp{peo} previously overlapped.
  \item Expected outcomes include duplication removal, contracting consolidation, and fewer stovepipes, with near-term authority friction during transition.
\end{itemize}

\subsection{Mapping: PAEs to PEOs and SYSCOMs}

\subsubsection{1) PAE Maritime}
\textbf{Scope:} Surface ships, combatants, logistics vessels, amphibious ships.

\textbf{Absorbing from SYSCOMs:}
\begin{itemize}
  \item Naval Sea Systems Command (primary).
\end{itemize}

\textbf{Absorbing from PEOs:}
\begin{itemize}
  \item \ac{peoships}.
  \item \ac{peoiws} (surface combat systems portion).
  \item \ac{peousc} (surface unmanned portion).
\end{itemize}

\textbf{What changes:}
\begin{itemize}
  \item Previously: hull (PEO Ships) and combat systems (IWS) split across authorities.
  \item Now: one authority over full ship system (hull plus combat system).
\end{itemize}

\textbf{Impact:} reduces the classic hull versus combat-system integration boundary.

\subsubsection{2) PAE Undersea}
\textbf{Scope:} Submarines, undersea warfare systems, and \acp{uuv}.

\textbf{Absorbing from SYSCOMs:}
\begin{itemize}
  \item Naval Sea Systems Command (submarine directorates).
\end{itemize}

\textbf{Absorbing from PEOs:}
\begin{itemize}
  \item Team Submarines elements (for platform portfolios).
  \item \ac{peouws}.
  \item \ac{peousc} (\ac{uuv} portion).
\end{itemize}

\textbf{What changes:}
\begin{itemize}
  \item Previously fragmented across platforms, payloads, and unmanned systems.
  \item Now unified into one undersea kill-chain authority.
\end{itemize}

\textbf{Impact:} enables end-to-end undersea portfolio management across platform, weapons, and unmanned systems.

\subsubsection{3) PAE Robotics and Autonomous Systems (RAS)}
\textbf{Scope:} Cross-domain autonomy (air, surface, and subsurface).

\textbf{Absorbing from SYSCOMs:}
\begin{itemize}
  \item Naval Air Systems Command.
  \item Naval Sea Systems Command.
  \item Space and Naval Warfare Systems Command / \ac{navwar} enterprise.
\end{itemize}

\textbf{Absorbing from PEOs:}
\begin{itemize}
  \item \ac{peousc}.
  \item Unmanned aviation elements aligned under NAVAIR-led portfolios.
  \item Autonomy software/control elements aligned under \ac{peoc4i}.
\end{itemize}

\textbf{What changes:}
\begin{itemize}
  \item Previously: unmanned capability split by domain.
  \item Now: technology-centric portfolio around autonomy stack plus platforms.
\end{itemize}

\textbf{Impact:} centralizes \ac{ai}/autonomy integration and weakens domain stovepipes from platform-centric acquisition.

\subsubsection{4) PAE Industrial Operations}
\textbf{Scope:} Shipyards, depots, sustainment, and logistics infrastructure.

\textbf{Absorbing from SYSCOMs:}
\begin{itemize}
  \item Naval Sea Systems Command (shipyards).
  \item Naval Air Systems Command (depots).
\end{itemize}

\textbf{Absorbing from other organizations:}
\begin{itemize}
  \item Naval Shipyards.
  \item Fleet Readiness Centers.
\end{itemize}

\textbf{What changes:}
\begin{itemize}
  \item Previously split among acquisition, sustainment, and industrial-base governance.
  \item Now unified as industrial capacity managed like a portfolio.
\end{itemize}

\textbf{Impact:} enables integrated control of maintenance backlog, shipyard throughput, and depot readiness with lifecycle cost/readiness optimization.

\subsubsection{5) PAE Strategic Systems Programs}
\textbf{Scope:} Nuclear deterrent systems (for example, \ac{ssbn}s and missiles).

\textbf{Absorbing from existing organization:}
\begin{itemize}
  \item Strategic Systems Programs.
\end{itemize}

\textbf{What changes:}
\begin{itemize}
  \item Minimal structural change because \ac{ssp} was already centralized.
  \item Now formally aligned inside the \ac{pae} construct.
\end{itemize}

\textbf{Impact:} primarily governance alignment rather than operational disruption.

\subsubsection{6) PAE Marine Corps}
\textbf{Scope:} \ac{usmc} ground systems and expeditionary capabilities.

\textbf{Absorbing from SYSCOMs:}
\begin{itemize}
  \item Marine Corps Systems Command.
\end{itemize}

\textbf{Absorbing from PEOs:}
\begin{itemize}
  \item PEO Land Systems (\ac{usmc}).
  \item \ac{peoc4i} (\ac{usmc} portions).
\end{itemize}

\textbf{What changes:}
\begin{itemize}
  \item Consolidates Marine Corps acquisition under one Navy-aligned \ac{pae}.
\end{itemize}

\textbf{Impact:} improves Navy-Marine Corps expeditionary alignment, with potential friction where Marine Corps requirements autonomy intersects portfolio controls.

\subsubsection{7) PAE Mission Systems (Source-Limited Update)}
\textbf{Scope:} \ac{c4i}, enterprise information technology, business systems, enterprise services, and mission-system integration across platform portfolios.

\textbf{Public source trail:}
\begin{itemize}
  \item On 16~Mar~2026, \ac{don} announced five additional \acp{pae} and stated that transition study efforts continued for aviation, industrial infrastructure, mission systems, and munitions, with those efforts to be announced as formally established~\autocite{DON-PAE-Establishment-2026}.
  \item DVIDS now has a public unit page titled ``Portfolio Acquisition Executive Mission Systems (PAE MS),'' but no associated public news stories or publications were posted on that unit page as of the access date~\autocite{DVIDS-PAE-MS-Unit-2026}.
  \item AFCEA's 18~Mar~2026 Navy Information Warfare Industry Day agenda listed Mr.\ James ``Jim'' Day, SES, Deputy Assistant Secretary of the Navy for Mission Systems, as the ``Portfolio Acquisition Executive'' speaker~\autocite{AFCEA-Navy-IW-Industry-Day-2026}.
  \item A public PEO C4I PMW~760 brief dated 22~Jan~2026 and marked Distribution Statement~A uses ``Integration in PAE Mission Systems,'' places PMW~760 ship \ac{c4i} integration in a PAE Mission Systems context, and shows explicit interfaces with PAE Maritime for \ac{scn} integration, software-defined warship, unmanned \ac{c4i}, \ac{mosa}/modularity, Aegis Guam, digital innovation, and \ac{aiml}~\autocite{PEOC4I-PMW760-NDIA-2026}.
\end{itemize}

\textbf{Likely legacy alignment:}
\begin{itemize}
  \item \ac{peoc4i}.
  \item \ac{peodigital}.
  \item \ac{peomlb}.
  \item Mission-system integration equities that cross platform \acp{peo}, especially ship \ac{c4i} integration with PAE Maritime.
\end{itemize}

\textbf{What changes:}
\begin{itemize}
  \item Previously: mission systems, ship integration, enterprise IT, business systems, and platform portfolios were coordinated across DASN Mission Systems, \ac{navwar}, PEO C4I, PEO Digital, PEO MLB, and platform \acp{peo}.
  \item Now: the board-prep answer should anticipate a single portfolio authority for mission-system delivery, with platform PAEs retaining platform outcomes and PAE Mission Systems driving integrated digital, C4I, and mission-system delivery.
\end{itemize}

\textbf{Impact:} reduces the separation between information-warfare acquisition, ship or aviation platform acquisition, and enterprise digital delivery. The exact transfer of billets, program offices, and technical authority remains an open public-source question until \ac{don} releases a formal charter, directive, or announcement.

\subsection{Cross-Cutting Changes (Most Important)}
\subsubsection{1) Collapse of PEO Boundaries}
\begin{quote}
PEO Ships builds hulls.\\
PEO IWS builds combat systems.\\
Different chains of authority.
\end{quote}

\begin{quote}
PAE Maritime owns both and can trade cost/performance internally.
\end{quote}

\textbf{Net effect:} lower inter-PEO coordination overhead.

\subsubsection{2) SYSCOM Role Shift}
\acp{syscom} are assessed to lose portions of decision authority while retaining engineering authority, technical standards, and infrastructure.

\textbf{Transition vector:} owner/operator toward service-provider functions.

\subsubsection{3) Portfolio-Level Trade Space}
\acp{pae} are assessed to gain latitude to shift funding between programs, optimize procurement and sustainment jointly, and manage industrial-base effects at portfolio level.

\textbf{Bottom line for reform:} the leverage is portfolio trade authority, not merely org-chart redraw.

\subsection{Risks and Friction Points}
\begin{description}[style=nextline, labelsep=0.5em, font=\bfseries]
  \item[1) Authority conflicts.] Tension between \ac{syscom} technical authority and \ac{pae} decision authority may drive early-cycle delays.
  \item[2) Workforce realignment.] \acp{peo} may lose decision scope and senior billets, creating talent migration risk.
  \item[3) Contracting consolidation.] Transition from program contracts to portfolio contracts may disrupt prime and supply-chain behavior.
\end{description}

\subsection{Bottom Line}
\begin{itemize}
  \item \acp{pae} absorb authority across both \acp{syscom} and \acp{peo}, not just one layer.
  \item The largest changes occur where multiple \acp{peo} historically split one system or one kill chain (Maritime, Undersea, RAS, and Mission Systems).
  \item The reform direction is portfolio-centric acquisition, closer to capability product lines.
\end{itemize}

\textit{Validation note: the existence of the six publicly announced baseline \acp{pae} is source-backed by the \ac{fy}~2026 \ac{ndaa} portfolio-acquisition-executive statute and the \ac{don} 16~Mar~2026 public announcement. PAE Mission Systems is supported here by public indicators and user-provided internal-announcement context, but a public \ac{don} establishment announcement or charter was not found during the 27~Apr~2026 source check. The detailed PEO/SYSCOM absorption mapping remains an analyst study aid until formal Navy implementation directives publish each portfolio boundary.}

\IfSubfilesClassLoaded{
  \RenewDocumentCommand{\entryname}{}{\textbf{\color{Modern} Acronym}}
  \RenewDocumentCommand{\descriptionname}{}{\textbf{\color{Modern} Definition}}
  \printnoidxglossary[
    type=\acronymtype,
    title=Acronyms,
    style=long-booktabs
  ]
}{}
\end{document}
